\label{hmt:p3:app:proofs}

Cet appendice rassemble les preuves qui peuvent être vérifiées sans reconstruire toute la narration. Les hypothèses apparaissent dans chaque énoncé. Les résultats classiques employés pour interpréter un opérateur déjà construit sont distingués des identités internes.

\section{Positivité de l’opérateur de vide}

\begin{lemma}
Si \(T=T^*\) et \(\Spec(T)\subset(0,1)\), alors \(I-T^n>0\) pour tout \(n\ge1\), \(I-T^{120}\) est inversible et
\[
\mathcal V=(I-T^{90})(I-T^{120})^{-1}>0.
\]
\end{lemma}

\begin{proof}
Pour \(t\in(0,1)\), \(1-t^n>0\). Par calcul fonctionnel, \(I-T^n\) est positif. Comme le spectre de \(I-T^{120}\) est séparé de zéro en dimension finie, il possède un inverse positif. Les deux facteurs sont des fonctions de \(T\), commutent et leur produit est positif.
\end{proof}

\begin{lemma}[Monotonie de la complétion]
La fonction
\[
r(q)=\frac{1-q^{90}}{1-q^{120}}
\]
est strictement décroissante sur \((0,1)\).
\end{lemma}

\begin{proof}
Avec \(s=q^{30}\),
\[
r(s)=\frac{1+s+s^2}{1+s+s^2+s^3}.
\]
Une dérivation donne
\[
\frac{dr}{ds}=
-\frac{s^2(3+2s+s^2)}{(1+s+s^2+s^3)^2}<0.
\]
Comme \(ds/dq=30q^{29}>0\), on a aussi \(dr/dq<0\).
\end{proof}

\section{Traces qui retrouvent les caractères constitutifs}

De \(\log\mathcal V=(\log r_+)P_++(\log r_-)P_-\) et \(\tau(P_\pm)=1/2\),

\[
2\tau(\log\mathcal V)=\log r_++\log r_-=\log(r_+r_-).
\]

De plus, \(RP_+=P_+\), \(RP_-=-P_-\), donc

\[
-2\tau(R\log\mathcal V)
=-\log r_++\log r_-
=\log(r_-/r_+).
\]

L’exponentiation produit \(\widehat\mu,\widehat\varepsilon,\widehat Z,\widehat c\). Il n’y a pas d’approximation à cette étape.

\section{Compatibilité avec la tour UHF}

Soit \(\mathcal A_m=M_2\otimes M_{9^m}\) et l’inclusion \(\iota_m(x)=x\otimes I_9\). L’espérance

\[
E_m=\mathrm{id}_{M_2}\otimes\mathrm{id}_{M_{9^m}}\otimes\tau_9
\]

satisfait \(E_m\iota_m=\mathrm{id}\). Par conséquent,

\[
E_m(\mathcal V\otimes I_{9^{m+1}})
=\mathcal V\otimes I_{9^m}.
\]

La même égalité vaut pour toute fonction continue de \(\mathcal V\), en particulier \(\log\mathcal V\). La trace normalisée est compatible et les caractères ne changent pas lors du raffinement.

\section{Coefficients du profil analytique de criticité}

Soit

\[
b_m=\frac{2(3m-1)}{\sqrt{899}}.
\]

Le développement \(1-\cos(b_mx)=b_m^2x^2/2-b_m^4x^4/24+O(x^6)\) donne

\[
u_0(x)=\frac1{24}\sum_{m=1}^{12}b_m^2x^2
-\frac1{288}\sum_{m=1}^{12}b_m^4x^4+O(x^6).
\]

Par conséquent,

\[
\sum b_m^2=\frac4{899}\sum(3m-1)^2=24,
\]

d’où le coefficient quadratique est égal à un. Le calcul entier du quatrième moment réduit le coefficient à

\[
-\frac{715769}{2424603}.
\]

Cette fraction apporte un contrôle algébrique plus informatif que sa valeur
décimale.

\section{Somme trigonométrique fermée}

L’identité

\[
\sum_{m=1}^{M}\cos(a+(m-1)d)
=\frac{\sin(Md/2)}{\sin(d/2)}
\cos\!\left(a+\frac{M-1}{2}d\right)
\]

avec \(M=12\), \(a=4y\), \(d=6y\) produit

\[
\sum_{m=1}^{12}\cos(2(3m-1)y)
=\frac{\sin36y}{\sin3y}\cos37y.
\]

La division par \(12\) prouve la première forme de \cref{mab:rm:eq:chaos-closed} ; \(2\sin36y\cos37y=\sin73y-\sin y\) prouve la seconde.

\section{Schwarziano negativo}

Pour \(x\in(0,1]\), chaque \(b_mx\in(0,\pi)\). Par conséquent

\[
u_0'(x)=\frac1{12}\sum b_m\sin(b_mx)>0,
\]

\[
u_0'''(x)=-\frac1{12}\sum b_m^3\sin(b_mx)<0.
\]

Comme \(f_\lambda'=-\lambda u_0'\), les facteurs \(-\lambda\) s’annulent dans les quotients de la dérivée schwarzienne, et

\[
S(f_\lambda)=\frac{u_0'''}{u_0'}
-\frac32(u_0''/u_0')^2<0.
\]

La parité étend la conclusion à \([-1,0)\). L’origine est un point critique quadratique d’après \(u_0(x)=x^2+O(x^4)\).

\section{Intégrale de Lévy}

La convergence dominée permet

\begin{align*}
-\int_0^1\log x\,d\mu_G(x)
&=\frac1{\log2}\int_0^1\frac{-\log x}{1+x}\,dx\\
&=\frac1{\log2}\sum_{n=0}^\infty(-1)^n
\int_0^1x^n(-\log x)\,dx\\
&=\frac1{\log2}\sum_{n=1}^\infty\frac{(-1)^{n-1}}{n^2}\\
&=\frac{\eta(2)}{\log2}
=\frac{\pi^2}{12\log2}.
\end{align*}

\section{Valeurs du caractère dodécaphasique}

Le produit \(\chi_{12}=\chi_{-3}\chi_{-4}\) s’annule si \(\gcd(n,12)>1\). Dans les quatre classes unités, la multiplication donne \((+,-,-,+)\). La série

\[
L(s,\chi_{12})=
\sum_{k\ge0}\left[
\frac1{(12k+1)^s}-\frac1{(12k+5)^s}
-\frac1{(12k+7)^s}+\frac1{(12k+11)^s}
\right]
\]

peut être évaluée en \(s=1\) au moyen de produits d’Euler ou de la formule des unités du corps \(\mathbb Q(\sqrt3)\), ce qui donne \cref{eq:L1chi12}. En \(s=2\), la formule de Bernoulli généralisée donne \cref{eq:L2chi12}. La multiplication directe du caractère constitue un critère de réfutation indépendant des deux évaluations.

\section{Reconstruction des secteurs aréal et modulaire}

Soit

\[
N=N_{90}+N_{120},\qquad D=90N_{90}+120N_{120}.
\]

La matrice du système est

\[
M=\begin{pmatrix}1&1\\90&120\end{pmatrix},
\qquad \det M=30.
\]

Par conséquent,

\[
M^{-1}=\frac1{30}
\begin{pmatrix}120&-1\\-90&1\end{pmatrix},
\]

et

\[
N_{90}=4N-D/30,\qquad N_{120}=D/30-3N.
\]

Le déterminant \(30\) explique pourquoi le pas modulaire fondamental réapparaît dans la classification infinie.

\section{Purification par un état thermochamp double}

Pour un qutrit de poids \(w>0\),

\[
Z(w)=1+\ee^{-w}+\ee^{-2w},
\qquad p_n=\frac{\ee^{-nw}}{Z(w)}.
\]

Définissez

\[
|\mathrm{TFD}(w)\rangle=\sum_{n=0}^2\sqrt{p_n}
|n\rangle_A|n\rangle_{\bar A}.
\]

Alors

\[
\Tr_{\bar A}|\mathrm{TFD}\rangle\langle\mathrm{TFD}|
=\sum_np_n|n\rangle\langle n|=\rho_A.
\]

L’égalité des entropies des deux sous-systèmes est immédiate par la décomposition de Schmidt. Le produit de dix-huit copies de chaque canal purifie l’état total de bord.

\section{1re loi modulaire aréale}

Pour une famille différentiable \(\rho(t)\), \(\Tr\rho(t)=1\), avec un état de référence \(\rho_0\) fidèle et \(K_0=-\log\rho_0\),

\[
\left.\frac{d}{dt}S(\rho(t))\right|_{t=0}
=\Tr(\dot\rho_0K_0),
\]

porque \(\Tr\dot\rho_0=0\). Si

\[
K_0-cI=\beta_{90}\mathcal A_{90}
+\beta_{120}\mathcal A_{120},
\]

le terme scalaire disparaît et

\[
\delta S=\beta_{90}\delta\langle\mathcal A_{90}\rangle
+\beta_{120}\delta\langle\mathcal A_{120}\rangle.
\]

Comme \(\beta_m=m\Dmod/(\gammaH a_0)\), le rapport est exactement \(120/90=4/3\).

\section{Réflexion rationnelle CKM}

Soit \(D=\operatorname{diag}(1,-1,1)\), \(\boldsymbol\theta=M\mathbf x\) et
\(R=MDM^{-1}\). Immédiatement, \(R^2=I\), \(\det R=\det D=-1\) et
\(RM=MD\). Dans les coordonnées déterminées,

\[
c_h=(-5,7,-20)^T,\qquad
\ell_h=(75,176,-150)/3857,\qquad \ell_hc_h=1.
\]

Le projecteur impair est \(c_h\ell_h\), de rang un. Ainsi

\[
R=I-2c_h\ell_h.
\]

La multiplication des fractions vérifie

\[
R=\frac1{3857}
\begin{pmatrix}
4607&1760&-1500\\
-1050&1393&2100\\
3000&7040&-2143
\end{pmatrix}.
\]

La phase \(\delta_{\rm CP}=9A\) reste fixe parce que \(D\) ne change pas \(A\). La conjugaison CP physique est un autre opérateur.

\section{Classification de la réalisation modulaire à la limite infinie}

Supposons une suite infinie de facteurs qutrit avec des états fidèles, dont les logarithmes des rapports de valeurs propres appartiennent à

\[
90\Dmod\Z\quad\text{ou}\quad120\Dmod\Z,
\]

et supposons que les deux types apparaissent avec une densité positive. Le groupe additif engendré est

\[
90\Dmod\Z+120\Dmod\Z=30\Dmod\Z.
\]

Le rapport asymptotique est alors

\[
\{0\}\cup\exp(-30\Dmod\Z)
=\{0\}\cup\lambda_\partial^\Z,
\qquad \lambda_\partial=\ee^{-30\Dmod}.
\]

Le théorème d’Araki--Woods pour les facteurs ITPFI classe la fermeture GNS comme le facteur hyperfini \(III_{\lambda_\partial}\). L’application est conditionnelle aux hypothèses de fidélité, de compatibilité et de persistance ; les poids qui interviennent dans la classification sont internes.

\begin{remark}
Si seul le canal \(90\) persiste, le groupe générateur devient \(90\Dmod\Z\) ; si seul le canal \(120\) persiste, il devient \(120\Dmod\Z\). Si les états locaux ne sont pas du type produit compatible déclaré, le critère doit être recalculé. Ces cas sont des critères de réfutation, non des notes de style.
\end{remark}


% REMATE_LOCAL_AP_J
\Needspace{22\baselineskip}
\section{Identité angulaire électrofaible}

Pour le caractère

\[
\chi(n_A,s_C)=\exp(n_AA+s_CC^*/6)
\]

et les routes \(W=(94,-1)\), \(Z=(95,-1)\),

\[
\frac{m_W}{m_Z}=\frac{\chi(94,-1)}{\chi(95,-1)}
=\ee^{-A}.
\]

Comme \(D_A=\ee^{-2A}=\det T\),

\[
(m_W/m_Z)^2=D_A,\qquad
\sin^2\theta_W^{\rm OS}=1-D_A.
\]

La preuve appartient à la carte angulaire. Son interprétation complète de
jauge exige l’entrelaceur de base déclaré dans le corps du volume.
